\documentclass[aspectratio=169,11pt]{beamer}
\usepackage[utf8]{inputenc}
\usepackage[T1]{fontenc}
\usepackage[spanish,es-noquoting]{babel}
\usepackage{tikz,booktabs}
\usepackage[absolute,overlay]{textpos}
\usetikzlibrary{shapes.geometric,arrows.meta,positioning,calc}

\usetheme{default}
% Sin usecolortheme: 'seabird' no viene en la instalacion estandar de
% TeX Live y los colores institucionales se definen abajo a mano.
\definecolor{itmazul}{RGB}{22,48,110}
\definecolor{acento}{RGB}{0,120,160}
\definecolor{verde}{RGB}{34,120,60}
\definecolor{rojo}{RGB}{170,40,40}
\setbeamercolor{structure}{fg=itmazul}
\setbeamercolor{frametitle}{fg=white,bg=itmazul}
\setbeamertemplate{navigation symbols}{}
\setbeamercolor{title}{fg=itmazul}
\setbeamercolor{itemize item}{fg=acento}
\setbeamercolor{itemize subitem}{fg=acento}
\setbeamercolor{block title}{fg=white,bg=itmazul}
\setbeamercolor{block body}{bg=itmazul!6}
\setbeamertemplate{footline}[frame number]
\setbeamerfont{frametitle}{size=\large}

\newcommand{\sect}[1]{\textcolor{acento}{\small\textbf{[Sec.~#1]}}}

% Marca de quien expone cada bloque, arriba a la derecha
\newcommand{\quien}[1]{%
  \begin{textblock*}{4cm}(8.6cm,0.32cm)
    \tiny\textcolor{white}{\textbf{#1}}
  \end{textblock*}}

\title{\textbf{Fan-out Masivo y Presencia en Tiempo Real}}
\subtitle{ComunicaMedellín — Canal de alertas ciudadanas}
\author{Grupo 3 \\[4pt]
  \small Julián \quad Elizabeth \quad Simón \quad Nilton \quad Lina}
\institute{Tópicos Avanzados de Bases de Datos · Módulo 1 · ITM}
\date{2026}

\begin{document}

\begin{frame}[plain]\titlepage\end{frame}

% ─────────────────────────────────────────────
\begin{frame}{El problema}
\quien{Julián}
\begin{columns}[T]
\begin{column}{0.55\textwidth}
\begin{itemize}
  \item \textbf{100.000+} suscriptores
  \item Una alerta debe llegar \textbf{casi simultáneamente}
  \item Guardar el mensaje en el historial de cada usuario
  \item Saber quién está \textbf{en línea} en todo momento
  \item Sin caerse ante un pico de tráfico
\end{itemize}
\vspace{6pt}
\footnotesize\textcolor{acento}{Cierres viales · alertas climáticas · cancelación de clases}
\end{column}
\begin{column}{0.42\textwidth}
\begin{tikzpicture}[scale=0.72,font=\scriptsize]
  \node[draw,rounded corners,fill=itmazul!10,minimum width=2.2cm] (a) at (0,2.4) {1 alerta};
  \foreach \i/\x in {1/-1.7,2/-0.55,3/0.6,4/1.75}
    \node[draw,circle,fill=acento!15,minimum size=0.5cm] (u\i) at (\x,0) {};
  \foreach \i in {1,2,3,4} \draw[-{Stealth[length=1.5mm]},acento] (a) -- (u\i);
  \node at (0,-0.75) {$\times$ 100.000 dispositivos};
\end{tikzpicture}
\end{column}
\end{columns}
\end{frame}

% ─────────────────────────────────────────────
\begin{frame}{Primero las cargas, después la tecnología}
\quien{Julián}
\centering\small
\begin{tabular}{@{}llll@{}}
\toprule
 & \textbf{Carga} & \textbf{Tipo} & \textbf{Sección} \\
\midrule
C1 & Publicar alerta & OLTP ligero & 2.1 \\
C2 & Fan-out a 100K & Streaming & 4.2 \\
C3 & Historial por usuario & OLTP escritura masiva & 1.3 \\
C4 & \textbf{Presencia} & \textbf{Clave-valor} & \textbf{1.2, 1.3} \\
\bottomrule
\end{tabular}
\vspace{12pt}
\begin{block}{Lo que descubrimos}
\small C2 y C4 \emph{parecen} lo mismo, pero exigen garantías opuestas:
la presencia tolera desactualización; la entrega de la alerta \textbf{no tolera pérdida}.\\
$\Rightarrow$ \textbf{Dos decisiones CAP dentro del mismo sistema.}
\end{block}
\end{frame}

% ─────────────────────────────────────────────
\begin{frame}{¿Por qué no basta un servidor?}
\quien{Julián}
\begin{columns}[T]
\begin{column}{0.48\textwidth}
\centering\textbf{Cliente-Servidor} \sect{3.2}
\begin{tikzpicture}[scale=0.68,font=\scriptsize]
  \node[draw,rounded corners,fill=rojo!12,minimum width=2cm,minimum height=0.8cm] (s) at (0,1.6) {servidor};
  \foreach \x in {-1.5,-0.5,0.5,1.5}
    \node[draw,circle,minimum size=0.35cm] at (\x,0) {};
  \foreach \x in {-1.5,-0.5,0.5,1.5} \draw[-{Stealth[length=1.2mm]},rojo] (s) -- (\x,0.25);
  \node[rojo] at (0,-0.7) {cuello de botella · SPOF};
\end{tikzpicture}
\end{column}
\begin{column}{0.48\textwidth}
\centering\textbf{Distribuida} \sect{3.3}
\begin{tikzpicture}[scale=0.68,font=\scriptsize]
  \foreach \i/\x in {1/-1.3,2/0,3/1.3}
    \node[draw,rounded corners,fill=verde!12,minimum width=1cm,minimum height=0.6cm] (n\i) at (\x,1.6) {n\i};
  \foreach \i/\x in {1/-1.3,2/0,3/1.3} {
    \node[draw,circle,minimum size=0.3cm] (u\i) at (\x,0) {};
    \draw[-{Stealth[length=1.2mm]},verde] (n\i) -- (u\i);}
  \node[verde] at (0,-0.7) {sharding + replicación};
\end{tikzpicture}
\end{column}
\end{columns}
\vspace{10pt}
\begin{alertblock}{}
\small Un PostgreSQL en la nube \textbf{sigue siendo} Cliente-Servidor.
«Estar en la nube» $\neq$ «ser distribuido» \sect{3.2 vs 3.3}.
\end{alertblock}
\end{frame}

% ─────────────────────────────────────────────
\begin{frame}{Lambda o Kappa: cómo se procesan los datos}
\quien{Elizabeth}
\begin{columns}[T]
\begin{column}{0.48\textwidth}
\textbf{Lambda} \sect{4.1}\\[2pt]
\footnotesize
Dos rutas de \emph{procesamiento}:
\begin{itemize}\itemsep1pt
\item capa \emph{batch}: recalcula el histórico
\item capa de velocidad: lo reciente
\item capa que fusiona ambas vistas
\end{itemize}
\vspace{3pt}
Precio: \textbf{dos bases de código} para la misma lógica.
\end{column}
\begin{column}{0.48\textwidth}
\textbf{Kappa} \sect{4.2} \quad\textcolor{verde}{\small\textbf{← la nuestra}}\\[2pt]
\footnotesize
Un solo \emph{pipeline} de \emph{streaming}:
\begin{itemize}\itemsep1pt
\item el log inmutable es la fuente de verdad
\item los consumidores derivan sus vistas
\item reprocesar = releer desde el \emph{offset} 0
\end{itemize}
\vspace{3pt}
Una sola lógica, latencia consistentemente baja.
\end{column}
\end{columns}

\vspace{6pt}
\begin{block}{\small El matiz que importa}
\footnotesize
Que haya \textbf{dos destinos} no lo convierte en Lambda. Lambda son dos formas
de \emph{procesar}. Nosotros tenemos \textbf{un pipeline con dos consumidores}:
no hay capa \emph{batch}, ni \emph{cron}, ni trabajos programados. Eso es Kappa,
y es además un sistema \textbf{orientado a eventos} \sect{4.3}.
\end{block}
\end{frame}

% ─────────────────────────────────────────────
\begin{frame}{Arquitectura}
\quien{Elizabeth}
\centering
\begin{tikzpicture}[font=\scriptsize,
  caja/.style={rounded corners=2pt,draw=itmazul!60,fill=itmazul!8,minimum height=0.75cm,align=center,inner sep=4pt},
  motor/.style={caja,fill=acento!14,draw=acento,thick},
  fl/.style={-{Stealth[length=2mm]},draw=itmazul!70}]
  \node[caja,minimum width=2cm] (api) {API};
  \node[motor,minimum width=7.4cm,below=0.55cm of api] (k)
    {\textbf{Kafka} \sect{4.2} — particiones: norte · centro · sur};
  \node[caja,below=0.55cm of k,xshift=-2.5cm,minimum width=1.5cm] (w1) {worker};
  \node[caja,below=0.55cm of k,minimum width=1.5cm] (w2) {worker};
  \node[caja,below=0.55cm of k,xshift=2.5cm,minimum width=1.5cm] (w3) {worker};
  \node[motor,below=0.6cm of w2,xshift=-2.1cm,minimum width=4cm] (c)
    {\textbf{Cassandra} ×3 · AP\\ \sect{1.3, 3.3}};
  \node[motor,below=0.6cm of w2,xshift=2.5cm,minimum width=3.2cm] (ch)
    {\textbf{ClickHouse}\\ \sect{2.4}};
  \draw[fl] (api) -- (k);
  \foreach \w in {w1,w2,w3} \draw[fl] (k.south) -- (\w.north);
  \draw[fl] (w1.south) -- (c.north west);
  \draw[fl] (w2.south) -- (c.north);
  \draw[fl,dashed,verde] (k.east) -- ++(0.6,0) |- (ch.east);
\end{tikzpicture}
\vspace{4pt}
\footnotesize Ruta operacional + ruta analítica sobre el mismo log $\Rightarrow$ Lambda \sect{4.1}
\end{frame}

% ─────────────────────────────────────────────
\begin{frame}{Decisión CAP}
\quien{Simón}
\begin{columns}[T]
\begin{column}{0.42\textwidth}
\begin{tikzpicture}[font=\scriptsize,scale=0.92]
  \def\R{1.75}
  \coordinate (C) at (90:\R); \coordinate (A) at (210:\R); \coordinate (P) at (330:\R);
  \fill[acento!8] (C)--(A)--(P)--cycle;
  \draw[thick,itmazul!70] (C)--(A)--(P)--cycle;
  \node[above] at (C) {\textbf{C}}; \node[below left] at (A) {\textbf{A}}; \node[below right] at (P) {\textbf{P}};
  \draw[very thick,rojo] ($(C)!0.5!(A)+(-0.12,0.12)$) -- ++(0.4,-0.4);
  \node[align=center] at (0,-0.3) {\textbf{AP}};
\end{tikzpicture}
\end{column}
\begin{column}{0.55\textwidth}
\textbf{Presencia $\rightarrow$ AP, nivel \texttt{ONE}}\\
\footnotesize Se sacrifica C. Un desfase de 2\,s en «en línea» \textbf{no cambia ninguna decisión}.\\[6pt]
\textbf{Historial $\rightarrow$ AP, escritura \texttt{QUORUM}}\\
\footnotesize La alerta no se puede perder: es seguridad pública.\\[8pt]
\scriptsize\textcolor{rojo}{Si eligiéramos CP:} durante una partición el sistema rechazaría heartbeats
y la app mostraría a todos desconectados — parecería caída.
\end{column}
\end{columns}
\end{frame}

% ─────────────────────────────────────────────
\begin{frame}{Contraste con TuBoleta}
\quien{Simón}
\centering\small
\begin{tabular}{@{}lll@{}}
\toprule
 & \textbf{TuBoleta} & \textbf{ComunicaMedellín} \\
\midrule
Recurso & Silla A-12: finito & Mensaje: replicable \\
Conflicto & Dos compran lo mismo & Nadie compite \\
Si falla C & Doble venta $\rightarrow$ demanda & Presencia 2\,s vieja \\
Decisión & \textcolor{rojo}{\textbf{CP}} (Raft) & \textcolor{verde}{\textbf{AP}} (ajustable) \\
Motor & CockroachDB & Cassandra \\
\bottomrule
\end{tabular}
\vspace{14pt}
\begin{block}{}
\small La consistencia fuerte cuesta latencia. Se paga cuando hay un
\textbf{recurso en disputa}. En presencia nadie «gana» el estar en línea:
pagar consenso sería gasto puro.
\end{block}
\end{frame}

% ─────────────────────────────────────────────
\begin{frame}{Por qué no otras soluciones}
\quien{Nilton}
\footnotesize
\begin{tabular}{@{}p{2.5cm}p{8.4cm}@{}}
\toprule
\textbf{Descartada} & \textbf{Razón} \\
\midrule
\textbf{CockroachDB} &
Es \textcolor{rojo}{\textbf{CP}} con consenso Raft: los nodos nunca divergen.
Correcta para TuBoleta, inservible aquí — sin niveles \texttt{ONE}/\texttt{QUORUM}
el reto sería irrealizable. Licencia BSL, no \emph{open source}. \sect{1.3, 3.3} \\
\addlinespace
\textbf{PostgreSQL} &
Cliente-Servidor: el servidor sigue siendo cuello de botella y punto único de
falla, esté en la nube o no. \sect{3.2} \\
\addlinespace
\textbf{MongoDB} &
Clasificada como \textcolor{rojo}{\textbf{CP}}. Prioriza consistencia sobre
disponibilidad: lo contrario de lo que exige un canal de alertas. \sect{1.2} \\
\addlinespace
\textbf{Redis} &
El módulo lo presenta como \emph{feature store} en memoria, no como base
distribuida. Sin consistencia ajustable por consulta. \sect{4.4} \\
\bottomrule
\end{tabular}

\vspace{5pt}
\begin{block}{\small}
\footnotesize
No descartamos CockroachDB por desconocerla, sino \textbf{porque sabemos para qué
sirve}: en un caso con recursos en disputa sería la elección correcta.
\end{block}
\end{frame}

% ─────────────────────────────────────────────
\begin{frame}{Demo en vivo}
\quien{Lina}
\footnotesize
\begin{enumerate}\itemsep2pt
  \item \texttt{make status} — 3 nodos \texttt{UN}, tokens repartidos, RF=3
  \item \textbf{Panel} — publicar una alerta: Kafka, 3 workers, Cassandra, dispositivos
  \item \textbf{Panel} — la fila \emph{recién creada} en \texttt{historial\_usuario}
  \item \texttt{make demo-consistencia} — \texttt{ONE} frente a \texttt{QUORUM}, 150 rondas
  \item \texttt{make demo-nodo-caido} — tumbar un nodo: \texttt{ALL} falla, y debe fallar
  \item \texttt{make ch-demo} — la agregación que Cassandra no puede resolver
\end{enumerate}

\vspace{4pt}
\begin{block}{\small Metodología}
\scriptsize
Los tres nodos comparten una máquina y replican en microsegundos: una condición
que \textbf{no existe en producción}. Con \texttt{tc netem} restituimos 250\,ms
\textbf{solo} al tráfico de replicación (puerto 7000), no a las lecturas del
cliente (9042). Sin eso, la demo mediría el laboratorio, no el sistema.
\end{block}
\end{frame}

% ─────────────────────────────────────────────
\begin{frame}{Resultados medidos}
\quien{Lina}
\begin{center}\small
\begin{tabular}{@{}lrrl@{}}
\toprule
\textbf{Nivel} & \textbf{Divergencias} & \textbf{P50} & \textbf{Lectura} \\
\midrule
\texttt{ONE} & \textcolor{rojo}{\textbf{149/150 (99{,}3\%)}} & \textbf{1{,}14 ms}
  & rápido, puede ir atrasado \\
\texttt{QUORUM} & \textcolor{verde}{\textbf{0/150 (0{,}0\%)}} & \textbf{130{,}3 ms}
  & consistente, 114$\times$ más lento \\
\bottomrule
\end{tabular}
\end{center}

\vspace{6pt}
\footnotesize
\begin{tabular}{@{}p{3.5cm}p{3.4cm}p{3.6cm}@{}}
\toprule
\textbf{Demostración} & \textbf{Medición} & \textbf{Qué prueba} \\
\midrule
Fan-out particionado & 10.127 entregas/s & el paralelismo es real \\
Reparto por zona & 34 / 33 / 33 & una zona = una partición \\
Caída de un nodo & \texttt{ONE} 3{,}1\,ms · \texttt{QUORUM} 752\,ms · \texttt{ALL} falla
  & a más consistencia, menos disponibilidad \\
Compresión columnar & 1\,M filas · \texttt{zona} 206:1 & por qué existe ClickHouse \\
\bottomrule
\end{tabular}

\vspace{4pt}
\centering\footnotesize
\textcolor{acento}{\textbf{Proyección: $\sim$10 s para los 100.000 suscriptores.}}
\end{frame}

% ─────────────────────────────────────────────
\begin{frame}{Extra: dos centros de datos}
\quien{Lina}
\footnotesize
Extendimos la topología al escenario de AeroMedellín: \texttt{cass1} y
\texttt{cass2} en el DC \textbf{medellin}, \texttt{cass3} en \textbf{miami},
con \texttt{NetworkTopologyStrategy \{medellin:2, miami:1\}} \sect{3.3}.

\vspace{6pt}
\begin{center}\small
\begin{tabular}{@{}lrl@{}}
\toprule
\textbf{Nivel} & \textbf{Lectura desde Miami} & \textbf{Qué hace} \\
\midrule
\texttt{ONE} & 1{,}58 ms & responde su réplica local \\
\texttt{LOCAL\_QUORUM} & \textcolor{verde}{\textbf{1{,}46 ms}} & mayoría \emph{dentro} de la ciudad \\
\texttt{QUORUM} & \textcolor{rojo}{\textbf{753{,}74 ms}} & mayoría global: cruza el océano \\
\texttt{EACH\_QUORUM} & 253{,}93 ms & mayoría en \emph{cada} ciudad \\
\bottomrule
\end{tabular}
\end{center}

\vspace{5pt}
\begin{block}{\small}
\footnotesize
\texttt{LOCAL\_QUORUM} da consistencia real por \textbf{516 veces menos latencia}
que \texttt{QUORUM} global. Por eso es el nivel estándar en producción multi-DC:
si Miami queda incomunicado, Medellín sigue operando con normalidad.
\end{block}
\end{frame}

% ─────────────────────────────────────────────
\begin{frame}{Conclusiones}
\begin{enumerate}\small
  \item Cuatro cargas $\Rightarrow$ \textbf{dos decisiones CAP} en un mismo sistema.
  \item La consistencia es un \textbf{parámetro por consulta}, no una propiedad global.
  \item Sacrificamos C \textbf{donde el impacto de negocio es nulo}, no por comodidad.
  \item Separar Cassandra (OLTP) de ClickHouse (OLAP) evita el conflicto de la \sect{2.2}.
\end{enumerate}
\vfill
\centering\footnotesize\texttt{github.com/lina082/comunicamedellin-grupo3}\\
\scriptsize Todo el stack es software libre (Apache 2.0 / MIT)
\end{frame}

\begin{frame}[plain]\centering\vfill
{\Large\textcolor{itmazul}{\textbf{¿Preguntas?}}}\\[8pt]
\small Equipo 3 — Módulo 1
\vfill\end{frame}

\end{document}
